\documentclass[../main.tex]{subfiles}
\begin{document}

Chương này trình bày cơ sở lý thuyết về hệ điều hành Windows, nhật ký giám sát Sysmon, và phân tích cấu trúc toán học của các mô hình phát hiện bất thường không giám sát (Unsupervised Anomaly Detection - UAD). Đồng thời, nghiên cứu giới thiệu các công cụ toán học tham số và phi tham số nhằm phục vụ việc phân tích phân phối dữ liệu trong không gian đa chiều.

\section{Nền tảng Windows Event Logging và nhật ký sự kiện Sysmon}
\label{section:1.1}

Hệ điều hành Windows sử dụng cơ chế ghi nhật ký sự kiện tập trung (\textbf{Windows Event Logging}) để ghi nhận các hoạt động của hệ thống dưới dạng tệp nhị phân \textit{.evtx}. Mỗi sự kiện bao gồm các thông tin cơ bản: thời gian xảy ra (\textit{Timestamp}), mã sự kiện (\textit{Event ID}), nguồn sinh nhật ký (\textit{Provider}), mức độ ưu tiên (\textit{Level}), và trường thông tin chi tiết (\textit{EventData}). 

Mặc dù Windows Event Log cung cấp khả năng giám sát hệ thống, các nhật ký mặc định thường thiếu thông tin chi tiết ở cấp tiến trình. Để phục vụ bài toán giám sát an ninh chuyên sâu, công cụ \textbf{Sysmon} (\textit{System Monitor}) thuộc bộ công cụ Microsoft Sysinternals được sử dụng như một lớp giám sát bổ sung có độ phân giải cao.

Sysmon cho phép ghi nhận chi tiết các hoạt động ở mức hệ thống. Các sự kiện từ Sysmon được sử dụng để xây dựng các đặc trưng hành vi tiến trình, bao gồm các Event ID được mô tả trong Bảng \ref{tab:sysmon_eids}:

\begin{table}[H]
\centering
\small
\setlength{\tabcolsep}{4pt}
\caption{Các Event ID chính của Sysmon phục vụ giám sát hành vi tiến trình}
\label{tab:sysmon_eids}
\begin{tabular}{|>{\centering\arraybackslash}p{1.8cm}|p{3.2cm}|p{8.0cm}|}
\hline
\textbf{Event ID} & \textbf{Tên Sự kiện} & \textbf{Ý nghĩa trong Giám sát An ninh} \\ \hline
\textbf{1} & \textit{Process Create} & Ghi nhận tiến trình được tạo lập, dòng lệnh (\textit{Command Line}), tiến trình cha. \\ \hline
\textbf{3} & \textit{Network Connection} & Ghi nhận kết nối mạng (IP, Port) liên kết trực tiếp với PID sinh ra nó. \\ \hline
\textbf{8} & \textit{CreateRemoteThread} & Giám sát hành vi can thiệp bộ nhớ tiến trình khác. \\ \hline
\textbf{10} & \textit{Process Access} & Ghi nhận thao tác mở tiến trình khác. \\ \hline
\textbf{11} & \textit{File Create} & Ghi nhận tệp tin mới được tạo dựng trên đĩa cứng. \\ \hline
\textbf{13} & \textit{Registry Value Set} & Theo dõi thay đổi cấu hình Registry. \\ \hline
\textbf{17} & \textit{Pipe Created} & Theo dõi giao tiếp liên tiến trình (IPC) qua Named Pipes. \\ \hline
\textbf{18} & \textit{Pipe Connected} & Theo dõi kết nối Named Pipe. \\ \hline
\textbf{22} & \textit{DNS Query} & Theo dõi yêu cầu phân giải tên miền sinh ra bởi tiến trình. \\ \hline
\end{tabular}
\end{table}

\subsection{Vấn đề tái sử dụng PID (PID Reuse)}
\label{subsection:1.1.1}
Trong hệ điều hành Windows, mỗi tiến trình khi khởi chạy được cấp một định danh số nguyên gọi là \textbf{Process ID (PID)}. Tuy nhiên, không gian PID bị giới hạn và Windows sẽ tái sử dụng (\textit{reuse}) các giá trị PID đã kết thúc cho các tiến trình mới khởi chạy sau đó. Điều này gây khó khăn trong việc liên kết các sự kiện log nếu chỉ sử dụng PID làm định danh duy nhất.

Để khắc phục vấn đề này, Sysmon cung cấp thuộc tính \textbf{ProcessGuid} (Global Unique Identifier) được tính toán dựa trên thời gian khởi tạo hệ thống và thời gian sinh tiến trình thực tế. ProcessGuid là định danh toàn cục không trùng lặp. Do đó, ProcessGuid được sử dụng làm khóa chính (\textit{primary key}) để định danh tiến trình trong dữ liệu.

\section{Biểu diễn tiến trình và chuẩn hóa đặc trưng (Không gian V10)}
\label{section:1.2}

Để áp dụng các mô hình học máy vào giám sát hệ thống, các sự kiện thô từ Sysmon cần được chuyển đổi thành các biểu diễn vector định lượng. Không gian đặc trưng V10 được định nghĩa là một không gian vector đa chiều đại diện cho hành vi của tiến trình.

\subsection{Chuẩn hóa đường dẫn tiến trình}
\label{subsection:1.2.1}
Trong hệ điều hành Windows, đường dẫn thực thi có thể xuất hiện dưới nhiều định dạng khác nhau (không phân biệt hoa thường, sử dụng biến môi trường). Toàn bộ đường dẫn tiến trình được chuẩn hóa qua các bước: chuyển về chữ thường (\textit{lowercase}), chuẩn hóa dấu phân tách thành \texttt{/}, và mở rộng toàn bộ các biến môi trường (\textit{environment variables expansion}) về dạng đường dẫn tuyệt đối.

\subsection{Biểu diễn đồ thị tiến trình}
\label{subsection:1.2.2}
Mối quan hệ thực thi giữa các tiến trình được mô hình hóa dưới dạng một đồ thị có hướng không chu trình (Directed Acyclic Graph - DAG): $G = (V, E)$. Trong đó:
\begin{itemize}
    \item Mỗi đỉnh $v \in V$ đại diện cho một vòng đời tiến trình duy nhất được định danh bằng \textbf{ProcessGuid}.
    \item Mỗi cạnh có hướng $e = (u, v) \in E$ đại diện cho quan hệ sinh tiến trình, được trích xuất từ sự kiện Sysmon Event ID 1.
\end{itemize}
Mỗi nút $v$ trên đồ thị chứa một tập hợp các thông số đặc trưng tĩnh và động trích xuất từ Event Data để xây dựng không gian đặc trưng hành vi (V10).

\section{Phân loại bất thường theo mô hình Chandola}
\label{section:1.3}
Theo nghiên cứu tổng quan của Chandola et al. (2009) \cite{chandola2009anomaly}, các hành vi bất thường trong hệ thống dữ liệu được chia làm ba nhóm chính:
\begin{itemize}
    \item \textbf{Bất thường điểm (Point Anomaly):} Một điểm dữ liệu đơn lẻ có giá trị khác biệt hoàn toàn so với phân phối bình thường.
    \item \textbf{Bất thường ngữ cảnh (Contextual Anomaly):} Điểm dữ liệu có thuộc tính hoàn toàn bình thường nếu đứng độc lập, nhưng trở thành bất thường khi đặt vào một ngữ cảnh cụ thể.
    \item \textbf{Bất thường tập thể (Collective Anomaly):} Tập hợp các sự kiện có cấu trúc bình thường khi xét riêng lẻ, nhưng sự xuất hiện đồng thời của chúng tạo nên một quy luật bất thường.
\end{itemize}

Sự phân loại này đóng vai trò quan trọng trong việc định hình yêu cầu đối với các hệ thống phát hiện bất thường: mô hình cần có khả năng phân tích đa chiều để bao quát cả ba loại bất thường trên.

\section{Tấn công có chủ đích (APT)}
\label{section:1.4}
Các chiến dịch tấn công có chủ đích (Advanced Persistent Threat - APT) thường liên quan đến các kỹ thuật xâm nhập phức tạp. Hệ thống giám sát an ninh mạng thường thu thập dữ liệu về các chiến dịch này dưới dạng các bộ dữ liệu (dataset). Trong nghiên cứu này, các thuật toán sẽ được đánh giá trên các tập dữ liệu mô phỏng lại một số chiến dịch APT đã biết, bao gồm:
\begin{itemize}
    \item Các kịch bản tấn công theo khuôn mẫu MITRE ATT\&CK (ví dụ: APT29).
    \item Các hành vi tạo kết nối mạng ẩn danh, truy xuất bộ nhớ, và sửa đổi cấu hình hệ thống (Registry).
\end{itemize}

\section{Lý thuyết Học máy không giám sát (UAD)}
\label{section:1.5}

Học máy không giám sát (Unsupervised Anomaly Detection - UAD) là phương pháp xây dựng mô hình hành vi bình thường (\textit{baseline}) từ dữ liệu không dán nhãn để phát hiện các cá thể bất thường. Các thuật toán UAD truyền thống dựa trên sự phân tách hình học hoặc cấu trúc phân bố. Dưới đây là lý thuyết của một số thuật toán điển hình:

\subsection{Isolation Forest (iForest)}
Isolation Forest \cite{liu2008isolation} dựa trên nguyên lý cô lập. Thuật toán chia cắt không gian dữ liệu bằng cách chọn ngẫu nhiên một đặc trưng và chọn ngẫu nhiên một điểm cắt nằm giữa giá trị lớn nhất và nhỏ nhất của đặc trưng đó. Quá trình chia cắt được lặp lại cho đến khi mọi điểm dữ liệu bị cô lập hoàn toàn dưới dạng các cây cô lập (Isolation Trees). Số lượng nhát cắt cần thiết để cô lập một điểm dữ liệu biểu thị mức độ bất thường của nó.

\subsection{Empirical Cumulative Distribution Functions (ECOD)}
ECOD là thuật toán phát hiện bất thường dựa trên xác suất đuôi thực nghiệm \cite{li2023ecod}. ECOD ước lượng Hàm phân phối tích lũy thực nghiệm (ECDF) cho từng chiều đặc trưng độc lập:
\begin{equation}
    F_{d}(x_d) = \frac{1}{n} \sum_{i=1}^{n} \mathbb{I}(X_{id} \le x_d)
\end{equation}
Điểm số rủi ro của điểm dữ liệu $x$ được tính bằng tổng log-likelihood xác suất đuôi của các chiều:
\begin{equation}
    Score(x) = -\sum_{d=1}^{D} \log\left( \min(F_{d}(x_d),\ 1 - F_{d}(x_d)) \right)
\end{equation}

\subsection{Local Outlier Factor (LOF)}
LOF là thuật toán phát hiện bất thường dựa trên mật độ cục bộ \cite{breunig2000lof}. LOF so sánh mật độ của một điểm dữ liệu với mật độ của $k$ láng giềng gần nhất của nó sử dụng khoảng cách Minkowski (thường là khoảng cách Euclidean). Một điểm có mật độ thấp hơn đáng kể so với láng giềng của nó sẽ được phân loại là bất thường.

\section{Khoảng cách toàn cục và ước lượng mật độ phi tham số}
\label{section:1.6}

Việc đo lường cấu trúc phân phối dữ liệu đa chiều thường dựa trên hai lý thuyết toán học: Khoảng cách toàn cục (Mahalanobis) và Ước lượng mật độ phi tham số (KDE).

\subsection{Khoảng cách Mahalanobis}
Khoảng cách Mahalanobis đo khoảng cách từ một điểm dữ liệu $x$ đến trung tâm phân phối $\mu$, có xét đến ma trận hiệp phương sai $\Sigma$ của toàn bộ dữ liệu:
\begin{equation}
    d_M(x) = \sqrt{(x - \mu)^T \Sigma^{-1} (x - \mu)}
\end{equation}
Khoảng cách này xấp xỉ phân phối của dữ liệu dưới dạng một ellipsoid toàn cục.

\subsection{Ước lượng mật độ phi tham số (KDE)}
Ước lượng mật độ nhân (Kernel Density Estimation - KDE) là phương pháp phi tham số để ước lượng hàm mật độ xác suất \cite{silverman1986density}:
\begin{equation}
    \hat{f}(x) = \frac{1}{nh} \sum_{i=1}^{n} K\left( \frac{x - X_i}{h} \right)
\end{equation}
Trong đó $h$ là tham số làm mịn (bandwidth). Đối với không gian dữ liệu hỗn hợp (chứa cả biến liên tục và rời rạc), mô hình \textbf{Mixed Product Kernel} của Racine và Li (2004) \cite{racine2004nonparametric} được áp dụng:
\begin{equation}
    \hat{f}(x) = \frac{1}{n} \sum_{i=1}^{n} \prod_{s=1}^{q} K_s(x_s, X_{is}, h_s)
\end{equation}
Với $K_s$ có thể là nhân Gaussian (cho biến liên tục), Aitchison-Aitken (cho biến phân loại), hoặc Wang-van Ryzin (cho biến phân loại có thứ tự).

\section{Khái niệm hòa trộn hành vi và định lượng đa chiều (BBS)}
\label{section:1.7}

Dựa trên sự phân loại bất thường của Chandola et al. (mục \ref{section:1.3}), nghiên cứu này đề xuất khái niệm \textbf{Không gian hòa trộn (Blending Space)} của mã độc. Kẻ tấn công có chủ đích (APT) thường áp dụng kỹ thuật \textit{Living off the Land} (LotL) để ngụy trang hành vi, trộn lẫn hoạt động của chúng vào hệ thống bình thường. Tuy nhiên, sự hòa trộn hiếm khi hoàn hảo ở mọi khía cạnh. Để đánh giá toàn diện, mức độ hòa trộn của mã độc được chiếu lên 3 trục độc lập (tương ứng với 3 loại bất thường):

\begin{itemize}
    \item \textbf{Sự hòa trộn Không gian (Spatial Blending - Trục 1):} Tương ứng với bất thường điểm. Mã độc cố gắng giả mạo các đặc trưng nội tại (ví dụ: tên tiến trình, tham số dòng lệnh) sao cho giống với tiến trình hệ thống nhất có thể trong không gian đặc trưng ($V_{10}$).
    \item \textbf{Sự hòa trộn Chuỗi thời gian (Sequential Blending - Trục 2):} Tương ứng với bất thường ngữ cảnh. Một tiến trình có vẻ bình thường nhưng thứ tự gọi thực thi của nó trong một khoảng thời gian lại bất thường so với thói quen của quản trị viên.
    \item \textbf{Sự hòa trộn Cấu trúc tập thể (Collective Blending - Trục 3):} Tương ứng với bất thường tập thể. Kẻ tấn công thường tạo ra hàng loạt tiến trình con với kiến trúc đồ thị liên kết bất thường (ví dụ: tiến trình mạng tự động sinh ra hàng loạt luồng shell nội bộ).
\end{itemize}

\subsection{Định lượng thông qua điểm Behavioral Blending Score (BBS)}
Khung đo lường BBS được thiết kế để chấm điểm sự hòa trộn trên 3 trục này với giá trị phân bố từ $0$ (lộ liễu hoàn toàn) đến $1$ (hòa trộn hoàn hảo). 

Theo quan điểm phòng thủ, hệ thống giám sát đa chiều sẽ bắt được kẻ tấn công nếu chúng để lộ sơ hở ở bất kỳ phương diện nào. Do đó, độ khó phát hiện chung của một tập dữ liệu (hay mức độ ngụy trang thực tế của mã độc) tuân theo \textbf{Nguyên lý Thùng gỗ (Min Aggregation)}. Điểm hòa trộn tổng quát ($BBS_{min}$) được quyết định bởi mặt có khả năng ngụy trang kém nhất:
\begin{equation}
    BBS_{min} = \min(BBS_{spa}, BBS_{seq}, BBS_{col})
\end{equation}
Cơ sở toán học để tính toán chi tiết các giá trị $BBS$ này được trình bày tại Chương 2.

\section{Các nghiên cứu liên quan}
\label{section:1.8}

Trong thập kỷ qua, các nghiên cứu nền tảng về chiều sâu dữ liệu (statistical data depth) \cite{liu1993quality, zuo2000general} đã chỉ ra đặc tính của các mô hình tham số (parametric models). Khoảng cách Mahalanobis, có thể được ước lượng thông qua các thuật toán như Minimum Covariance Determinant (MCD) \cite{rousseeuw1999fast}, cung cấp ước lượng cấu trúc toàn cục mạnh mẽ. 

Việc xây dựng các bộ dữ liệu chuẩn (benchmarks) để đánh giá hệ thống phát hiện bất thường đã được cộng đồng nghiên cứu quan tâm. Các nghiên cứu như của Janssens \& Postma (2012) \cite{janssens2012data} hay Zimek et al. (2012) \cite{zimek2012survey} phân tích tác động của số chiều dữ liệu. Công trình của Emmott et al. (2013) \cite{emmott2013systematic} và Campos et al. (2016) \cite{campos2016evaluation} cho thấy các thước đo cấu trúc phân phối có ảnh hưởng lớn đến kết quả thuật toán.

Các nghiên cứu của Davis \& Goadrich (2006) \cite{davis2006relationship} và Saito \& Rehmsmeier (2015) \cite{saito2015precision} đã phân tích tính chất của đường cong ROC trên dữ liệu mất cân bằng. Sự ra đời của các phương pháp phát hiện cục bộ (Local Outlier Detection) như LOCI \cite{papadimitriou2003loci} hay LODA \cite{pevny2016loda} nhằm bổ sung cho phân phối toàn cục. 

\end{document}
